\documentclass[11pt]{article}
\usepackage{amsmath,amssymb}
\usepackage[margin=1in]{geometry}

\newcommand{\bra}[1]{\langle #1|}
\newcommand{\ket}[1]{|#1\rangle}
\newcommand{\anti}[4]{\langle #1#2\Vert #3#4\rangle}
\newcommand{\dir}[4]{\langle #1#2|#3#4\rangle}
\newcommand{\cc}{^{*}}

\title{Unrestricted Coupled Cluster Singles and Doubles\\ for Complex Orbitals}
\author{}
\date{}

\begin{document}
\maketitle

\section{Scope}

The standard spin-orbital CCSD equations (Stanton, Gauss, Watts, Bartlett,
J.\ Chem.\ Phys.\ \textbf{94}, 4334 (1991)) are written for real orbitals. In
that form several integrals appear in a particular index order that is only
correct because of the real-orbital symmetry $\anti{p}{q}{r}{s}=\anti{r}{s}{p}{q}$.
Here the equations are re-derived by tracking which operator string every
integral multiplies, so that they hold for \emph{arbitrary complex}
spin-orbitals $\varphi_p(\mathbf r,\sigma)$. This covers

\begin{itemize}
\item UCCSD with complex spatial orbitals, $\varphi_p=\psi_p(\mathbf r)\,\sigma_p$ with
      $\sigma_p\in\{\alpha,\beta\}$ (Section~\ref{sec:blocks});
\item general two-component or four-component spinors (the case needed for the
      relativistic code), where there is no spin factorization at all.
\end{itemize}

No time-reversal or other symmetry is assumed. The reference determinant
$\ket{0}$ is built from the $N$ occupied spin-orbitals and need not be a
Hartree--Fock solution (the Fock matrix may have non-zero occupied--virtual
blocks).

\section{Conventions}

Indices $i,j,k,l,m,n$ label occupied, $a,b,c,d,e,f$ virtual, $p,q,r,s$ general
spin-orbitals. The one-electron operator is $h$, with
$h_{pq}=\bra{p}h\ket{q}=\int\varphi_p^*\,h\,\varphi_q$. Integrals:
\begin{align}
\dir{p}{q}{r}{s} &= \iint \varphi_p^*(1)\,\varphi_q^*(2)\,\frac{1}{r_{12}}\,\varphi_r(1)\,\varphi_s(2),\\
\anti{p}{q}{r}{s} &= \dir{p}{q}{r}{s}-\dir{p}{q}{s}{r}.
\end{align}
The Hamiltonian in second quantization,
\begin{equation}
H=\sum_{pq}h_{pq}\,a_p^\dagger a_q+\tfrac14\sum_{pqrs}\anti{p}{q}{r}{s}\,
a_p^\dagger a_q^\dagger a_s a_r ,
\end{equation}
so that $\anti{p}{q}{r}{s}$ is always the coefficient of the string
$a_p^\dagger a_q^\dagger a_s a_r$: \emph{the two left indices are the created
orbitals, the two right indices the annihilated ones.} This single rule fixes
every conjugation below.

The Fock matrix and the reference energy are
\begin{equation}
f_{pq}=h_{pq}+\sum_i\anti{p}{i}{q}{i},\qquad
E_0=\sum_i h_{ii}+\tfrac12\sum_{ij}\anti{i}{j}{i}{j},
\end{equation}
$f_{pq}=f_{qp}^*$ (Hermitian, in general complex).

\subsection*{Integral symmetries}

\begin{center}
\begin{tabular}{lcc}
 & real orbitals & complex orbitals\\ \hline
$\anti{p}{q}{r}{s}=-\anti{q}{p}{r}{s}=-\anti{p}{q}{s}{r}=\anti{q}{p}{s}{r}$ & yes & yes\\
$\anti{p}{q}{r}{s}=\anti{r}{s}{p}{q}$ & yes & \textbf{no}\\
$\anti{p}{q}{r}{s}=\anti{r}{s}{p}{q}^{*}$ & (trivial) & yes\\
$f_{pq}=f_{qp}$ & yes & \textbf{no}, $f_{pq}=f_{qp}^*$\\
\end{tabular}
\end{center}
The 8-fold permutational symmetry of the real case is reduced to 4-fold for the
unantisymmetrized $\dir{p}{q}{r}{s}$ (namely
$\dir{p}{q}{r}{s}=\dir{q}{p}{s}{r}=\dir{r}{s}{p}{q}^*=\dir{s}{r}{q}{p}^*$).
In particular, the blocks with virtuals in the bra and occupied in the ket
(e.g.\ $\anti{a}{b}{i}{j}$) are the complex conjugates of the corresponding
blocks with occupied in the bra (e.g.\ $\anti{i}{j}{a}{b}$), and are not equal
to them.

\section{Cluster ansatz and amplitude equations}

\begin{equation}
\ket{\Psi}=e^{T}\ket{0},\qquad T=T_1+T_2,
\end{equation}
\begin{equation}
T_1=\sum_{ia}t_i^a\,a_a^\dagger a_i,\qquad
T_2=\tfrac14\sum_{ijab}t_{ij}^{ab}\,a_a^\dagger a_b^\dagger a_j a_i ,
\end{equation}
with $t_{ij}^{ab}$ antisymmetric in $(ij)$ and in $(ab)$. The amplitudes are
complex and satisfy no conjugation relation such as
$t_{ij}^{ab}=t_{ji}^{ba*}$. The equations are the usual projections of the
similarity-transformed Hamiltonian $\bar H=e^{-T}He^{T}$,
\begin{equation}
E_c=\bra{0}\bar H\ket{0}-E_0,\qquad
\bra{\Phi_i^a}\bar H\ket{0}=0,\qquad
\bra{\Phi_{ij}^{ab}}\bar H\ket{0}=0 .
\end{equation}
Neither the $T$ amplitudes nor the left-hand projectors are ever complex
conjugated. Complex orbitals therefore enter \emph{only} through the
integrals, and the algebra is identical to the real case term by term; what
changes is which of the two members of a conjugate pair
$\anti{p}{q}{r}{s}$, $\anti{r}{s}{p}{q}$ appears. There is no left-hand
state and no $\Lambda$ equation in what follows.

\section{Working equations}

\subsection{Correlation energy}

\begin{equation}
E_c=\sum_{ia}f_{ia}\,t_i^a+\tfrac14\sum_{ijab}\anti{i}{j}{a}{b}\,t_{ij}^{ab}
+\tfrac12\sum_{ijab}\anti{i}{j}{a}{b}\,t_i^a t_j^b .
\end{equation}
For the exact wave function $E_c$ is real. For truncated CCSD no identity
guarantees it, so take $\operatorname{Re}E_c$ and monitor
$\operatorname{Im}E_c$ (PySCF's GCCSD warns when it exceeds $10^{-4}$).

\subsection{Intermediates}

With
\begin{equation}
\tilde\tau_{ij}^{ab}=t_{ij}^{ab}+\tfrac12\left(t_i^a t_j^b-t_i^b t_j^a\right),\qquad
\tau_{ij}^{ab}=t_{ij}^{ab}+t_i^a t_j^b-t_i^b t_j^a ,
\end{equation}
and $P(pq)X_{pq}=X_{pq}-X_{qp}$:
\begin{align}
F_{ae} &= (1-\delta_{ae})f_{ae}-\tfrac12\sum_m f_{me}\,t_m^a
 +\sum_{mf}t_m^f\anti{m}{a}{f}{e}
 -\tfrac12\sum_{mnf}\tilde\tau_{mn}^{af}\anti{m}{n}{e}{f},\\
F_{mi} &= (1-\delta_{mi})f_{mi}+\tfrac12\sum_e t_i^e f_{me}
 +\sum_{ne}t_n^e\anti{m}{n}{i}{e}
 +\tfrac12\sum_{nef}\tilde\tau_{in}^{ef}\anti{m}{n}{e}{f},\\
F_{me} &= f_{me}+\sum_{nf}t_n^f\anti{m}{n}{e}{f},\\
W_{mnij} &= \anti{m}{n}{i}{j}+P(ij)\sum_e t_j^e\anti{m}{n}{i}{e}
 +\tfrac14\sum_{ef}\tau_{ij}^{ef}\anti{m}{n}{e}{f},\\
W_{abef} &= \anti{a}{b}{e}{f}-P(ab)\sum_m t_m^b\anti{a}{m}{e}{f}
 +\tfrac14\sum_{mn}\tau_{mn}^{ab}\anti{m}{n}{e}{f},\\
W_{mbej} &= \anti{m}{b}{e}{j}+\sum_f t_j^f\anti{m}{b}{e}{f}
 -\sum_n t_n^b\anti{m}{n}{e}{j}
 -\sum_{nf}\left(\tfrac12 t_{jn}^{fb}+t_j^f t_n^b\right)\anti{m}{n}{e}{f}.
\end{align}
Each of these multiplies, in $\bar H$, an operator string whose created
indices are the left indices (e.g.\ $F_{ae}$ multiplies $a_a^\dagger a_e$ and
$W_{abef}$ multiplies $a_a^\dagger a_b^\dagger a_f a_e$), and every
integral inside them keeps the created-left / annihilated-right pattern of
that string. This is why these seven definitions need no conjugate: all
integrals appearing there are already in the orientation that the operator
algebra delivers.

\subsection{$T_1$ equation}

\begin{multline}
t_i^a D_i^a = f_{ai}+\sum_e t_i^e F_{ae}-\sum_m t_m^a F_{mi}
+\sum_{me}t_{im}^{ae}F_{me}
-\sum_{nf}t_n^f\anti{n}{a}{i}{f}\\
-\tfrac12\sum_{mef}t_{im}^{ef}\anti{m}{a}{e}{f}
-\tfrac12\sum_{men}t_{mn}^{ae}\anti{n}{m}{e}{i}.
\end{multline}

\subsection{$T_2$ equation}

\begin{multline}
t_{ij}^{ab}D_{ij}^{ab}=\anti{a}{b}{i}{j}
+P(ab)\sum_e t_{ij}^{ae}\Big(F_{be}-\tfrac12\sum_m t_m^b F_{me}\Big)
-P(ij)\sum_m t_{im}^{ab}\Big(F_{mj}+\tfrac12\sum_e t_j^e F_{me}\Big)\\
+\tfrac12\sum_{mn}\tau_{mn}^{ab}W_{mnij}
+\tfrac12\sum_{ef}\tau_{ij}^{ef}W_{abef}
+P(ij)P(ab)\sum_{me}\Big(t_{im}^{ae}W_{mbej}-t_i^e t_m^a\anti{m}{b}{e}{j}\Big)\\
+P(ij)\sum_e t_i^e\anti{a}{b}{e}{j}
-P(ab)\sum_m t_m^a\anti{m}{b}{i}{j}.
\end{multline}

The denominators contain the diagonal of the (Hermitian, hence real-diagonal)
Fock matrix,
\begin{equation}
D_i^a=f_{ii}-f_{aa},\qquad D_{ij}^{ab}=f_{ii}+f_{jj}-f_{aa}-f_{bb}.
\end{equation}
The off-diagonal Fock elements enter through the $(1-\delta)$ terms in
$F_{ae},F_{mi}$, so no canonical-orbital assumption is made. Level shifting
and DIIS act on the complex amplitude vector as in the real case, with the
DIIS overlap matrix taken as $\mathbf e_k^\dagger\mathbf e_l$ and the
extrapolation coefficients real.

\section{Differences with respect to the real-orbital equations}

Compared with the equations as usually printed for real orbitals, only the two
inhomogeneous terms change, and they change because the real-orbital
formulas are written with the occupied-bra orientation using
$\anti{p}{q}{r}{s}=\anti{r}{s}{p}{q}$:

\begin{center}
\begin{tabular}{lll}
term & real-orbital form & complex form (this note)\\ \hline
$T_1$ inhomogeneous term & $f_{ia}$ & $f_{ai}=f_{ia}^*$\\
$T_2$ inhomogeneous term & $\anti{i}{j}{a}{b}$ & $\anti{a}{b}{i}{j}=\anti{i}{j}{a}{b}^*$
\end{tabular}
\end{center}

All other terms, including $P(ij)\,t_i^e\anti{a}{b}{e}{j}$ and
$-P(ab)\,t_m^a\anti{m}{b}{i}{j}$, are printed with the same index order as in
the real case; they are correct for complex orbitals because that order is the
created-left / annihilated-right order of the operator algebra. What differs
is only how they are obtained from stored arrays: if the code stores the
occupied-bra blocks ($\anti{i}{j}{m}{b}$, $\anti{j}{e}{b}{a}$), then
\begin{equation}
\anti{a}{b}{e}{j}=\anti{j}{e}{b}{a}^*,\qquad
\anti{m}{b}{i}{j}=\anti{i}{j}{m}{b}^*,
\end{equation}
(pair swap $\anti{p}{q}{r}{s}=\anti{q}{p}{s}{r}$ combined with Hermitian
conjugation) and the conjugate must be taken. For real orbitals every
conjugation is trivial and the equations reduce to Stanton's. The same
pattern is used in PySCF's \texttt{GCCSD} (\texttt{pyscf/cc/gccsd.py},
\texttt{update\_amps}), which can serve as an external reference with complex
spin-orbitals.

\section{Initial guess}

With $f_{ov}=0$ the first iteration from $T=0$ gives the MP2 amplitudes
\begin{equation}
t_i^a=\frac{f_{ai}}{D_i^a},\qquad
t_{ij}^{ab}=\frac{\anti{a}{b}{i}{j}}{D_{ij}^{ab}},
\end{equation}
and $E_{\mathrm{MP2}}=\tfrac14\sum\anti{i}{j}{a}{b}\,t_{ij}^{ab}
=\tfrac14\sum|\anti{i}{j}{a}{b}|^2/D_{ij}^{ab}$ is real and negative as it
should be. Using $\anti{i}{j}{a}{b}$ instead of $\anti{a}{b}{i}{j}$ in the
numerator would give a complex energy, which is a convenient sanity check on
the orientation of the integrals.

\section{UCCSD: spin blocks for complex spatial orbitals}
\label{sec:blocks}

For $\varphi_p=\psi_p\,\sigma_p$ with complex spatial $\psi_p$,
\begin{equation}
\dir{p}{q}{r}{s}=\delta_{\sigma_p\sigma_r}\delta_{\sigma_q\sigma_s}\,
(pr|qs),\qquad
(pr|qs)=\iint\psi_p^*(1)\psi_r(1)\,\frac1{r_{12}}\,\psi_q^*(2)\psi_s(2),
\end{equation}
and $h$, $f$ conserve spin. Spin conservation then restricts the amplitudes
to
\begin{itemize}
\item $t_i^a$: $\sigma_i=\sigma_a$, giving $t^{\alpha}$ and $t^{\beta}$;
\item $t_{ij}^{ab}$: $\sigma_i+\sigma_j=\sigma_a+\sigma_b$, giving the
      unique blocks $\alpha\alpha\to\alpha\alpha$, $\beta\beta\to\beta\beta$,
      and $\alpha\beta\to\alpha\beta$ ($t_{iJ}^{aB}$; the blocks with $a,B$
      swapped follow by antisymmetry).
\end{itemize}
The equations of the previous section hold verbatim in this block structure:
the sums run over the spin-allowed combinations and the integral blocks are
$\anti{p}{q}{r}{s}=(pr|qs)-(ps|qr)$ for equal spins (all four orbitals of one spin, or
$\sigma_p=\sigma_r,\sigma_q=\sigma_s$) and $\anti{p}{q}{r}{s}=(pr|qs)$ for the
mixed block $p,r\in\alpha$, $q,s\in\beta$. The conjugation pattern of the
previous section applies unchanged to each block; in particular the
$\alpha\beta$ inhomogeneous term is
$\anti{a}{B}{i}{J}=(ai|BJ)=(ia|JB)^*$, with no exchange part.

For genuine spinors (two- or four-component) $\varphi_p$ is not an eigenfunction
of $\sigma_z$, the integrals $\dir{p}{q}{r}{s}$ do not factorize, and the
spin-orbital equations of Section~4 are used directly with all index
combinations, i.e.\ as a generalized CCSD (GCCSD) with complex orbitals.

\section{Cost and storage}

The scaling is that of real CCSD, $O(o^2v^4)$ time and $O(o^2v^2)$ amplitude
storage, with complex arithmetic (about $4\times$ the flops and $2\times$ the
memory per element, using \texttt{zgemm} in the contractions). Because
$\anti{p}{q}{r}{s}$ has only 4-fold symmetry, one cannot store only the unique
integrals of one orientation and reuse them for the conjugate orientation by
a transposition; the conjugate blocks of the previous section are obtained
by conjugating the stored block, never by transposing it alone.

\section{Left-hand equations ($\Lambda$) and reduced density matrices}
\label{sec:lambda}

Everything above is the right-hand side of the CCSD problem. Densities, and any property that is not the energy, need
the left-hand state $\bra{0}(1+\Lambda)e^{-T}$,
\begin{equation}
\Lambda=\sum_{ia}\lambda_i^a\,a_i^\dagger a_a+\tfrac14\sum_{ijab}\lambda_{ij}^{ab}\,a_i^\dagger a_j^\dagger a_b a_a ,
\end{equation}
with $\lambda_{ij}^{ab}$ antisymmetric in $(ij)$ and in $(ab)$. The $\Lambda$ equations are the stationarity conditions of
the Lagrangian
\begin{equation}
\mathcal L(t,\lambda)=E_c(t)+\sum_{ia}\lambda_i^a\,R_i^a(t)+\tfrac14\sum_{ijab}\lambda_{ij}^{ab}\,R_{ij}^{ab}(t),
\qquad R_\mu(t)=\bra{\Phi_\mu}\bar H\ket{0}=n_\mu(t)-D_\mu t_\mu ,
\end{equation}
where $n_\mu$ is the numerator of Section~4 ($t_\mu=n_\mu/D_\mu$ is the amplitude equation). $\mathcal L$ is holomorphic in
$t$ and $\lambda$ (neither is ever complex conjugated), so $\partial\mathcal L/\partial t_\mu=0$ gives \emph{linear}
equations for $\lambda$ with the same complex integrals as the $T$ equations; at convergence $\mathcal L=E_c$. With complex
orbitals the only change with respect to the real-orbital formulas is again which member of a conjugate pair of integrals
appears. The inhomogeneous terms are $f_{ia}$ and $\anti{i}{j}{a}{b}$ with the occupied indices in the bra and
\emph{not} conjugated (the complex conjugates of the $T$-equation driving terms); the conjugated integrals
$\overline{\anti{j}{k}{i}{c}}=\anti{i}{c}{j}{k}$ and $\overline{\anti{j}{a}{c}{b}}=\anti{b}{c}{a}{j}$ enter only through
$W^{\rm ovoo}$ and $W^{\rm vvvo}$ below. The equations are those of PySCF's \texttt{gccsd\_lambda.py}
(\texttt{make\_intermediates}, \texttt{update\_lambda}) in the conventions of this note, written out from the same data that
was evaluated with \texttt{numpy} and compared with PySCF (Section~\ref{sec:lambda-check}).

\subsection{$T$-dependent intermediates}

$\tau_{ij}^{ab}=t_{ij}^{ab}+t_i^at_j^b-t_i^bt_j^a$ as in Section~4. The intermediate $F_{kc}$ below equals $F_{me}$ of
Section~4.2. $\mathcal A,\mathcal B$ are one-body, $\mathcal C,\mathcal D,\mathcal E,\mathcal G$ and the four $W$ two-body
building blocks; all are independent of $\lambda$ and built once.
\begin{align*}
  \mathcal A_{ba} &=  f_{ba} -  \sum_{j}\,f_{ja}\,t_{j}^{b} \notag\\
  & \quad -  \sum_{jc}\,\anti{j}{b}{a}{c}\,t_{j}^{c} + \tfrac{1}{2}\, \sum_{jkc}\,\anti{j}{k}{c}{a}\,\tau_{jk}^{bc}
\end{align*}
\begin{align*}
  \mathcal B_{ij} &=  f_{ij} +  \sum_{b}\,f_{ib}\,t_{j}^{b} \notag\\
  & \quad -  \sum_{kb}\,\anti{k}{i}{j}{b}\,t_{k}^{b} + \tfrac{1}{2}\, \sum_{kbc}\,\anti{i}{k}{b}{c}\,\tau_{jk}^{bc}
\end{align*}
\begin{align*}
  \mathcal C_{ijkl} &=  \sum_{cd}\,\anti{i}{j}{c}{d}\,\tau_{kl}^{cd}
\end{align*}
\begin{align*}
  \mathcal D_{jcbk} &=  \sum_{ld}\,\anti{l}{j}{d}{b}\,t_{kl}^{cd} +  \anti{j}{c}{b}{k}
\end{align*}
\begin{align*}
  F_{kc} &=  f_{kc} -  \sum_{ld}\,\anti{k}{l}{d}{c}\,t_{l}^{d}
\end{align*}
\begin{align*}
  \mathcal E_{bj} &=  f_{bj} +  \sum_{kc}\,f_{kc}\,t_{jk}^{bc} \notag\\
  & \quad +  \sum_{kc}\,F_{kc}\,t_{k}^{b}\,t_{j}^{c} - \tfrac{1}{2}\, \sum_{klc}\,\anti{k}{l}{j}{c}\,t_{kl}^{bc} \notag\\
  & \quad + \tfrac{1}{2}\, \sum_{kcd}\,\anti{k}{b}{d}{c}\,t_{jk}^{cd}
\end{align*}
\begin{align*}
  \mathcal G_{ck} &=  \mathcal E_{ck} +  \sum_{jb}\,\mathcal D_{jcbk}\,t_{j}^{b} \notag\\
  & \quad +  \sum_{b}\,\mathcal A_{cb}\,t_{k}^{b} -  \sum_{j}\,\mathcal B_{jk}\,t_{j}^{c}
\end{align*}
\begin{align*}
  W^{\rm oooo}_{ijkl} &= \tfrac{1}{2}\, \anti{i}{j}{k}{l} + \tfrac{1}{4}\, \mathcal C_{ijkl} \notag\\
  & \quad +  \sum_{c}\,\anti{i}{j}{k}{c}\,t_{l}^{c}
\end{align*}
\begin{align*}
  W^{\rm ovvo}_{jcbk} &=  \mathcal D_{jcbk} -  \sum_{ld}\,\anti{l}{j}{d}{b}\,t_{l}^{c}\,t_{k}^{d} \notag\\
  & \quad -  \sum_{l}\,\anti{l}{j}{k}{b}\,t_{l}^{c} +  \sum_{d}\,\anti{j}{c}{b}{d}\,t_{k}^{d}
\end{align*}
\begin{align*}
  W^{\rm ovoo}_{icjk} &= \tfrac{1}{4}\, \sum_{bd}\,\anti{i}{c}{d}{b}\,\tau_{jk}^{db} + \tfrac{1}{2}\, \overline{\anti{j}{k}{i}{c}} \notag\\
  & \quad +  \sum_{b}\,\mathcal D_{icbk}\,t_{j}^{b} -  \sum_{lb}\,\anti{l}{i}{j}{b}\,t_{kl}^{cb}
\end{align*}
\begin{align*}
  W^{\rm vvvo}_{bcak} &=  \sum_{j}\,\mathcal D_{jcak}\,t_{j}^{b} + \tfrac{1}{4}\, \sum_{jl}\,\anti{j}{l}{k}{a}\,\tau_{jl}^{bc} \notag\\
  & \quad - \tfrac{1}{2}\, \overline{\anti{k}{a}{c}{b}} +  \sum_{jd}\,\anti{j}{b}{a}{d}\,t_{kj}^{cd}
\end{align*}

\subsection{$\Lambda$ update}

The $\lambda$-dependent intermediates are
\begin{align*}
  M^{\rm v}_{ba} &= \tfrac{1}{2}\, \sum_{klc}\,\lambda_{kl}^{ca}\,t_{kl}^{cb}
\end{align*}
\begin{align*}
  M^{\rm o}_{ij} &= \tfrac{1}{2}\, \sum_{kcd}\,\lambda_{ki}^{cd}\,t_{kj}^{cd}
\end{align*}
\begin{align*}
  N^{\rm v}_{ba} &=  M^{\rm v}_{ba} +  \sum_{k}\,\lambda_{k}^{a}\,t_{k}^{b}
\end{align*}
\begin{align*}
  N^{\rm o}_{ik} &=  M^{\rm o}_{ik} +  \sum_{c}\,\lambda_{i}^{c}\,t_{k}^{c}
\end{align*}
\begin{align*}
  S_{ijkl} &=  \sum_{cd}\,\lambda_{ij}^{cd}\,\tau_{kl}^{cd}
\end{align*}
\begin{align*}
  U_{ijck} &=  \sum_{d}\,\lambda_{ij}^{cd}\,t_{k}^{d}
\end{align*}
\begin{align*}
  M_{ijab} &=  \sum_{kl}\,\lambda_{kl}^{ab}\,W^{\rm oooo}_{ijkl} + \tfrac{1}{4}\, \sum_{kl}\,\anti{k}{l}{a}{b}\,S_{ijkl} \notag\\
  & \quad -  \sum_{kc}\,\anti{k}{c}{b}{a}\,U_{ijck} + \tfrac{1}{2}\, \sum_{cd}\,\lambda_{ij}^{cd}\,\anti{c}{d}{a}{b}
\end{align*}
\begin{align*}
  q_{jb} &=  t_{j}^{b} +  \sum_{kc}\,\lambda_{k}^{c}\,t_{kj}^{cb} \notag\\
  & \quad -  \sum_{d}\,N^{\rm v}_{bd}\,t_{j}^{d} -  \sum_{l}\,M^{\rm o}_{lj}\,t_{l}^{b}
\end{align*}

The term $\tfrac12\sum_{cd}\lambda_{ij}^{cd}\anti{c}{d}{a}{b}$ in $M_{ijab}$ is the $\Lambda$-equation particle--particle ladder: like the $T$-equation ladder it needs only the $vvvv$ slabs, and $\sum_{c}\anti{k}{c}{b}{a}\,U_{ijck}$ needs the $ovvv$ slabs, so the streamed scheme of Section~\ref{sec:impl} applies unchanged.
The diagonal Fock elements are moved to the denominators, $\mathcal A'_{ba}=\mathcal A_{ba}-\delta_{ba}f_{aa}$ and
$\mathcal B'_{ij}=\mathcal B_{ij}-\delta_{ij}f_{ii}$ (they are not part of the right-hand sides). With
$P(ij)X_{ijab}=X_{ijab}-X_{jiab}$ and $P(ab)X_{ijab}=X_{ijab}-X_{ijba}$ the numerator of the doubles equation is
\begin{equation}
\mathcal N_{ij}^{ab}=\anti{i}{j}{a}{b}+M_{ijab}+P(ij)P(ab)\,X_{ijab}-P(ab)\,Y_{ijab}+P(ij)\,Z_{ijab},
\end{equation}
with
\begin{align*}
  X_{ijab} &=  \lambda_{i}^{a}\,F_{jb} +  \sum_{kc}\,\lambda_{ki}^{ca}\,W^{\rm ovvo}_{jcbk}
\end{align*}
\begin{align*}
  Y_{ijab} &=  \sum_{k}\,\lambda_{k}^{a}\,\anti{i}{j}{k}{b} +  \sum_{c}\,\lambda_{ij}^{ca}\,\mathcal A'_{cb} \notag\\
  & \quad +  \sum_{c}\,N^{\rm v}_{ca}\,\anti{i}{j}{c}{b}
\end{align*}
\begin{align*}
  Z_{ijab} &=  \sum_{c}\,\lambda_{j}^{c}\,\anti{i}{c}{a}{b} +  \sum_{k}\,\lambda_{ki}^{ab}\,\mathcal B'_{jk} \notag\\
  & \quad -  \sum_{k}\,N^{\rm o}_{ik}\,\anti{k}{j}{a}{b}
\end{align*}
and that of the singles equation
\begin{align*}
  \mathcal N_{i}^{a} &=  f_{ia} +  \sum_{jb}\,\lambda_{j}^{b}\,\anti{i}{b}{a}{j} \notag\\
  & \quad +  \sum_{b}\,\lambda_{i}^{b}\,\mathcal A'_{ba} -  \sum_{j}\,\lambda_{j}^{a}\,\mathcal B'_{ij} \notag\\
  & \quad -  \sum_{jkc}\,\lambda_{kj}^{ca}\,W^{\rm ovoo}_{icjk} -  \sum_{kbc}\,\lambda_{ik}^{bc}\,W^{\rm vvvo}_{bcak} \notag\\
  & \quad +  \sum_{jb}\,M_{ijab}\,t_{j}^{b} +  \sum_{jb}\,\lambda_{ji}^{ba}\,\mathcal G_{bj} \notag\\
  & \quad +  \sum_{jb}\,\anti{j}{i}{b}{a}\,q_{jb} +  \sum_{bc}\,\anti{i}{c}{a}{b}\,N^{\rm v}_{bc} \notag\\
  & \quad -  \sum_{jk}\,\anti{j}{i}{k}{a}\,N^{\rm o}_{kj} -  \sum_{k}\,M^{\rm o}_{ik}\,F_{ka} \notag\\
  & \quad -  \sum_{c}\,M^{\rm v}_{ca}\,F_{ic}
\end{align*}
Both are divided by the same denominators as the $T$ equations,
\begin{equation}
\lambda_i^a\leftarrow\frac{\mathcal N_i^a}{D_i^a},\qquad
\lambda_{ij}^{ab}\leftarrow\frac{\mathcal N_{ij}^{ab}}{D_{ij}^{ab}},\qquad
D_i^a=f_{ii}-f_{aa},\quad D_{ij}^{ab}=f_{ii}+f_{jj}-f_{aa}-f_{bb},
\end{equation}
iterated from $\lambda=t$ with DIIS and an optional level shift exactly like the amplitudes. The cost per iteration is that
of one $T$ iteration (the same contractions, with $\lambda$ replacing $t$ in the ladder).

\subsection{One-particle density matrix}

Let $\gamma_{pq}=\bra{0}(1+\Lambda)e^{-T}a_p^\dagger a_q e^{T}\ket{0}$, so that the one-body energy is
$\sum_{pq}h_{pq}\gamma_{pq}$ (orbital relaxation is not included). In the occupied--virtual partitioning
$\gamma_{pq}=\delta_{pq}n_p+D_{qp}$, with the blocks $D^{\rm oo}_{ij}$, $D^{\rm vv}_{ab}$, $D^{\rm vo}_{ai}$ and
$D^{\rm ov}_{ia}=\lambda_i^a$ given by
\begin{align*}
  D^{\rm oo}_{ij} &= -  \sum_{e}\,\lambda_{i}^{e}\,t_{j}^{e} - \tfrac{1}{2}\, \sum_{mef}\,\lambda_{im}^{ef}\,t_{jm}^{ef}
\end{align*}
\begin{align*}
  D^{\rm vv}_{ab} &=  \sum_{m}\,t_{m}^{a}\,\lambda_{m}^{b} + \tfrac{1}{2}\, \sum_{mne}\,t_{mn}^{ea}\,\lambda_{mn}^{eb}
\end{align*}
\begin{align*}
  x^{(1)}_{mi} &= \tfrac{1}{2}\, \sum_{nef}\,\lambda_{mn}^{ef}\,t_{in}^{ef}
\end{align*}
\begin{align*}
  x^{(2)}_{ae} &= \tfrac{1}{2}\, \sum_{mnf}\,t_{mn}^{fa}\,\lambda_{mn}^{fe} +  \sum_{m}\,t_{m}^{a}\,\lambda_{m}^{e}
\end{align*}
\begin{align*}
  D^{\rm vo}_{ai} &=  \sum_{me}\,t_{im}^{ae}\,\lambda_{m}^{e} -  \sum_{m}\,x^{(1)}_{mi}\,t_{m}^{a} \notag\\
  & \quad -  \sum_{e}\,t_{i}^{e}\,x^{(2)}_{ae} +  t_{i}^{a}
\end{align*}
The coupled-cluster (Lagrangian) density is not Hermitian, already for real orbitals (e.g.\ $D^{\rm ov}$ and
$D^{\rm vo}$ differ); the Hermitian part, which is what the real part of the energy and of its derivatives
involves, is
\begin{equation}
\bar\gamma_{pq}=\tfrac12\big(\gamma_{pq}+\gamma_{qp}^{*}\big),\qquad
\operatorname{Re}\sum_{pq}h_{pq}\gamma_{pq}=\sum_{pq}h_{pq}\bar\gamma_{pq}.
\end{equation}
PySCF returns $\bar\gamma^{\mathsf T}$ (its \texttt{dm1[p,q]} is $\bar\gamma_{qp}=\langle a_q^\dagger a_p\rangle$).

\subsection{Two-particle density matrix}

Let $\Gamma_{pqrs}=\bra{0}(1+\Lambda)e^{-T}a_p^\dagger a_q^\dagger a_s a_r e^{T}\ket{0}$, so the energy is
\begin{equation}
E=\sum_{pq}h_{pq}\bar\gamma_{pq}+\tfrac14\sum_{pqrs}\anti{p}{q}{r}{s}\,\Gamma_{pqrs}\quad(\text{real part}),
\end{equation}
and $\Gamma$ is antisymmetric in $(pq)$ and in $(rs)$, with $\Gamma_{pqrs}^{*}=\Gamma_{rspq}$ and
$\sum_p\Gamma_{pqps}=(N-1)\bar\gamma_{qs}$. It is assembled in four steps; the notation is that of \texttt{gccsd\_rdm.py}.
In this part (and only here) the intermediate $\bar\tau_{ij}^{ab}=t_{ij}^{ab}+2t_i^at_j^b$ is \emph{not} antisymmetrized:
the antisymmetrization is carried out by the $d$-blocks below, and using the antisymmetrized $\tau$ in these formulas
changes the result.

\paragraph{Intermediates.}
\begin{align*}
  \bar\tau_{ij}^{ab} &=  t_{ij}^{ab} + 2\, t_{i}^{a}\,t_{j}^{b}
\end{align*}
\begin{align*}
  m_{iajb} &=  \sum_{kc}\,\lambda_{ik}^{ac}\,t_{kj}^{cb}
\end{align*}
\begin{align*}
  u^{(1)}_{ia} &=  \sum_{kc}\,\lambda_{k}^{c}\,t_{ki}^{ca}
\end{align*}
\begin{align*}
  u^{(2)}_{cb} &=  \sum_{k}\,\lambda_{k}^{c}\,t_{k}^{b}
\end{align*}
\begin{align*}
  u^{(3)}_{kj} &=  \sum_{c}\,\lambda_{k}^{c}\,t_{j}^{c}
\end{align*}
\begin{align*}
  u^{(4)}_{lj} &=  \sum_{d}\,m_{ldjd}
\end{align*}
\begin{align*}
  u^{(5)}_{db} &=  \sum_{l}\,m_{ldlb}
\end{align*}
\begin{align*}
  u^{(6)}_{ijkl} &= \tfrac{1}{16}\, \sum_{cd}\,\lambda_{kl}^{cd}\,\bar\tau_{ij}^{cd}
\end{align*}
\begin{align*}
  u^{(7)}_{ij} &= \tfrac{1}{4}\, \sum_{c}\,m_{icjc}
\end{align*}
\begin{align*}
  u^{(8)}_{ab} &= \tfrac{1}{4}\, \sum_{k}\,m_{kakb}
\end{align*}

\paragraph{$g$-blocks.} The $oovv$, $ovvv$ and $ooov$ blocks are complex conjugated before the last (unconjugated) term of
the last two is added:
\begin{align*}
  g^{\rm oovv}_{ijab} &= \Big[\tfrac{1}{4}\, \overline{\lambda_{ij}^{ab}} + \tfrac{1}{4}\, \bar\tau_{ij}^{ab} \notag\\
  & \quad +  u^{(1)}_{ia}\,t_{j}^{b} + \tfrac{1}{2}\, \sum_{c}\,u^{(2)}_{cb}\,t_{ij}^{ca} \notag\\
  & \quad + \tfrac{1}{2}\, \sum_{k}\,\bar\tau_{ki}^{ab}\,u^{(3)}_{kj} - \tfrac{1}{4}\, \sum_{l}\,u^{(4)}_{lj}\,\bar\tau_{li}^{ba} \notag\\
  & \quad - \tfrac{1}{4}\, \sum_{d}\,u^{(5)}_{db}\,\bar\tau_{ji}^{da} - \tfrac{1}{2}\, \sum_{ld}\,m_{ldia}\,\bar\tau_{lj}^{bd} \notag\\
  & \quad +  \sum_{kl}\,u^{(6)}_{ijkl}\,\bar\tau_{kl}^{ab}\Big]^{*}
\end{align*}
\begin{align*}
  g^{\rm vvvv}_{abcd} &= \tfrac{1}{8}\, \sum_{ij}\,\bar\tau_{ij}^{ab}\,\lambda_{ij}^{cd}
\end{align*}
\begin{align*}
  g^{\rm oooo}_{klij} &= \tfrac{1}{8}\, \sum_{ab}\,\lambda_{kl}^{ab}\,\bar\tau_{ij}^{ab}
\end{align*}
\begin{align*}
  g^{\rm ooov}_{jkia} &= \Big[- \tfrac{1}{4}\, \sum_{b}\,\bar\tau_{jk}^{ba}\,\lambda_{i}^{b} +  \sum_{l}\,g^{\rm oooo}_{iljk}\,t_{l}^{a} \notag\\
  & \quad -  u^{(7)}_{ij}\,t_{k}^{a} + \tfrac{1}{2}\, \sum_{c}\,m_{icja}\,t_{k}^{c}\Big]^{*}\ + \tfrac{1}{4}\, \sum_{b}\,\lambda_{jk}^{ab}\,t_{i}^{b}
\end{align*}
\begin{align*}
  g^{\rm ovvo}_{ibaj} &=  \lambda_{i}^{a}\,t_{j}^{b} +  m_{iajb} \notag\\
  & \quad -  \sum_{kc}\,\lambda_{ik}^{ac}\,t_{j}^{c}\,t_{k}^{b}
\end{align*}
\begin{align*}
  g^{\rm ovvv}_{iabc} &= \Big[\tfrac{1}{4}\, \sum_{j}\,\lambda_{j}^{a}\,\bar\tau_{ij}^{bc} +  \sum_{d}\,g^{\rm vvvv}_{bcad}\,t_{i}^{d} \notag\\
  & \quad +  u^{(8)}_{ac}\,t_{i}^{b} + \tfrac{1}{2}\, \sum_{k}\,m_{kaib}\,t_{k}^{c}\Big]^{*}\ + \tfrac{1}{4}\, \sum_{j}\,\lambda_{ij}^{bc}\,t_{j}^{a}
\end{align*}
Here $g^{\rm oooo}$ in $g^{\rm ooov}$ is the unconjugated block, and the $g^{\rm ovvo}$ and $g^{\rm vvvv}$ blocks
enter the $d$-blocks without conjugation.

\paragraph{$d$-blocks.} Antisymmetrizing and symmetrizing the $g$-blocks,
\begin{align*}
  d^{\rm ovov}_{iajb} &=  g^{\rm oovv}_{ijab} - g^{\rm oovv}_{ijba} \\
  d^{\rm vvvv}_{abcd} &=  g^{\rm vvvv}_{acbd} - g^{\rm vvvv}_{acdb} \\
  d^{\rm oooo}_{ijkl} &=  g^{\rm oooo}_{ikjl} - g^{\rm oooo}_{iklj} \\
  d^{\rm ovvv}_{iabc} &=  g^{\rm ovvv}_{ibac} - g^{\rm ovvv}_{ibca} \\
  d^{\rm ooov}_{ijka} &=  g^{\rm ooov}_{ikja} - g^{\rm ooov}_{kija} \\
  d^{\rm ovvo}_{iabj} &=  g^{\rm ovvo}_{ibaj}
\end{align*}
and then
\begin{align*}
 d^{\rm ovov}_{iajb}&\leftarrow\tfrac12\big(d^{\rm ovov}_{iajb}+d^{\rm ovov}_{jbia}\big),\\
 d^{\rm vvvv}_{abcd}&\leftarrow d^{\rm vvvv}_{abcd}+\overline{d^{\rm vvvv}_{badc}},\qquad
 d^{\rm oooo}_{ijkl}\leftarrow d^{\rm oooo}_{ijkl}+\overline{d^{\rm oooo}_{jilk}},\\
 d^{\rm ovvo}_{iabj}&\leftarrow\tfrac12\big(d^{\rm ovvo}_{iabj}+\overline{d^{\rm ovvo}_{jbai}}\big).
\end{align*}

\paragraph{Assembly.} Define the array $\mathcal R_{pqrs}\equiv\Gamma_{qspr}$ blockwise ($i,j,k$ occupied; $a,b,c$ virtual),
\begin{align*}
  \mathcal R_{iajb} &= d^{\rm ovov}_{iajb} \\
  \mathcal R_{aibj} &= \overline{d^{\rm ovov}_{iajb}} \\
  \mathcal R_{ijab} &= -d^{\rm ovvo}_{ibaj} \\
  \mathcal R_{abij} &= -d^{\rm ovvo}_{ibaj} \\
  \mathcal R_{iabj} &= d^{\rm ovvo}_{iabj} \\
  \mathcal R_{aijb} &= \overline{d^{\rm ovvo}_{iabj}} \\
  \mathcal R_{abcd} &= d^{\rm vvvv}_{abcd} \\
  \mathcal R_{ijkl} &= d^{\rm oooo}_{ijkl} \\
  \mathcal R_{iabc} &= d^{\rm ovvv}_{iabc} \\
  \mathcal R_{abic} &= d^{\rm ovvv}_{icab} \\
  \mathcal R_{abci} &= \overline{d^{\rm ovvv}_{icba}} \\
  \mathcal R_{aibc} &= \overline{d^{\rm ovvv}_{iacb}} \\
  \mathcal R_{ijka} &= d^{\rm ooov}_{ijka} \\
  \mathcal R_{iajk} &= d^{\rm ooov}_{jkia} \\
  \mathcal R_{ijak} &= \overline{d^{\rm ooov}_{jika}} \\
  \mathcal R_{aijk} &= \overline{d^{\rm ooov}_{kjia}}
\end{align*}
where the block of $\mathcal R$ is selected by the occupied/virtual character of $p,q,r,s$ in that order.
The correlation part of the one-particle matrix, $\bar\delta_{pq}\equiv\bar\gamma_{qp}-\delta_{pq}n_p$ (with $n_p=1$ for
occupied $p$), is then added for every occupied $i$,
\begin{equation}
\mathcal R_{ii\,pq}\mathrel{+}=\bar\delta_{pq},\quad \mathcal R_{pq\,ii}\mathrel{+}=\bar\delta_{pq},\quad
\mathcal R_{p\,ii\,q}\mathrel{-}=\bar\delta_{pq},\quad \mathcal R_{i\,pq\,i}\mathrel{-}=\bar\delta_{qp},
\end{equation}
followed by the mean-field part, $\mathcal R_{iijj}\mathrel{+}=1$ and $\mathcal R_{ijji}\mathrel{-}=1$ for all occupied
$i,j$, and finally
\begin{equation}
\Gamma_{pqrs}=\mathcal R_{rpsq}.
\end{equation}
(For real orbitals every conjugation is trivial and $\Gamma$ reduces to the usual 2-RDM of spin-orbital CCSD.)

\subsection{Storage}

Of the $g$/$d$-blocks only $g^{\rm vvvv}$ is $O(v^4)$. To avoid ever storing it, contract it directly with the
integrals: $\tfrac14\sum\anti{a}{b}{c}{d}\,\Gamma_{abcd}$ with $\Gamma_{abcd}=d^{\rm vvvv}_{cadb}$ and
$g^{\rm vvvv}_{abcd}=\tfrac18\sum_{ij}\bar\tau_{ij}^{ab}\lambda_{ij}^{cd}$ is evaluated slab by slab exactly like the
$T$ ladder, $O(o^2v^4)$ time and no $v^4$ memory. Any other use of $\Gamma_{abcd}$ needs it explicitly.

\subsection{Density matrix files (\texttt{CC\_RDM})}

\texttt{src/CC/lambda\_Uccsd.\{h,cpp\}} solves the $\Lambda$ equations above (Jacobi sweeps with DIIS from $\lambda=t$,
the same convergence criterion as $T$; the $vvvv$ ladder is streamed one slab $\langle c\,d|a\,b\rangle$ at a time, the
$ovvv$ block $\anti{m}{a}{e}{f}$ is held once) and writes
\begin{itemize}
\item \texttt{<input>\_MO.dm1}: records $(\texttt{int32}\;p,q;\ \texttt{float64}\;\operatorname{Re},\operatorname{Im})$ with
 $\bar\gamma_{pq}=\langle a_p^\dagger a_q\rangle$ (Hermitian part), all $p,q$;
\item \texttt{<input>\_MO.dm2}: records $(\texttt{int32}\;p,q,r,s;\ \texttt{float64}\;\operatorname{Re},\operatorname{Im})$ with
 $\Gamma_{pqrs}=\langle a_p^\dagger a_q^\dagger a_s a_r\rangle$ for $p<q$, $r<s$ only (the remaining elements follow from
 $\Gamma_{pqrs}=-\Gamma_{qprs}=-\Gamma_{pqsr}$); the energy is
 $E=\sum_{pq}h_{pq}\bar\gamma_{pq}+\tfrac12\sum_{pqrs}\langle pq|rs\rangle\Gamma_{pqrs}$;
\end{itemize}
both ended by an all-zero record. Indices are 1-based and in Kramers-pair order, $(1,2)=(\psi_1,\bar\psi_1)$,
$(3,4)=(\psi_2,\bar\psi_2),\dots$: X2C and the positive-energy block of the 4-component run are numbered that way
already, the non-relativistic spin-orbitals are renumbered $\alpha_k\to2k+1$, $\beta_k\to2k+2$. Frozen-core spinors are
included with their mean-field contributions (they appear only through the one-body and mean-field terms of
Section~9.4). Elements with $|\Gamma|$ below \texttt{CC\_RDM\_THRESHOLD} are omitted, and the energy recomputed from the
written elements is printed next to $E_0+E_c$ as a check (the two agree to $10^{-9}$ Eh for the default $10^{-10}$).
The 2-RDM is produced element by element from the $d$-blocks (never as a dense $n^4$ array); only the all-virtual
block is formed slab by slab from $\bar\tau$ and $\lambda$. The densities are unrelaxed: no orbital response is included.

\subsection{Densities in the basis of the Gaussian primitives (\texttt{src/CC/gauss\_rdm.\{h,cpp\}})}

With \texttt{CC\_RDM\_GAUSSIAN} (default) the same densities are also written over the Cartesian Gaussian \emph{primitives}
$\chi_\mu=N_\mu x^{l_x}y^{l_y}z^{l_z}e^{-\alpha_\mu|\mathbf r-\mathbf A|^2}$ in the formats read by \texttt{RHO2\_OPS}:
\texttt{<input>\_GAU.basis}, \texttt{<input>\_GAU.dm1}, \texttt{<input>\_GAU.dm2}.

\paragraph{Expansion.} Every contracted Cartesian AO is unrolled into its primitives (the libcint factor of its shell times its
contraction coefficient is the weight of the plain monomial), and identical primitives (same centre, exponent and powers) are
merged. A spinor has one coefficient vector per spin/large--small component $c$: $\varphi_p^{(c)}=\sum_\mu C^{c}_{\mu p}\chi_\mu$,
with $c\in\{\alpha,\beta\}$ over the Large primitives for NON\_REL and X2C (the NON\_REL spin-orbitals are pure $\alpha$ or $\beta$,
the spherical transformation and the \texttt{LOWGEN} reduction are folded into $C$), and $c\in\{L\alpha,L\beta,S\alpha,S\beta\}$ for
the 4-component spinors, the Small component being expanded over the raw RKB small AOs (their own primitives, kept distinct from the
Large ones). The weights are divided by $N_\mu=(2\alpha/\pi)^{3/4}(4\alpha)^{l/2}/\sqrt{(2l_x-1)!!(2l_y-1)!!(2l_z-1)!!}$ (the
normalization \texttt{RHO2\_OPS} multiplies by); \texttt{CC\_RDM\_CNORM1 TRUE} keeps it inside $D$ instead (\texttt{\$CNORM1}).

\paragraph{Definition.} The spin- and component-summed pair density is
$\rho_2(\mathbf r_1,\mathbf r_2)=\sum D_{\mu\nu\lambda\sigma}\,\chi_\mu(\mathbf r_1)\chi_\lambda(\mathbf r_1)\chi_\nu(\mathbf r_2)\chi_\sigma(\mathbf r_2)$,
\begin{equation}
D^{\rm raw}_{\mu\nu\lambda\sigma}=\sum_{c_1c_2}\sum_{pqrs}\Gamma_{pqrs}\;
\overline{C^{c_1}_{\mu p}}\,C^{c_1}_{\lambda r}\;\overline{C^{c_2}_{\nu q}}\,C^{c_2}_{\sigma s},
\end{equation}
so only pairs of primitives of the same component occur ($\mu,\lambda$ both Large or both Small, and likewise
$\nu,\sigma$), and $\sum D_{\mu\nu\lambda\sigma}S_{\mu\lambda}S_{\nu\sigma}=N(N-1)$. The real basis functions make
$\chi_\mu\chi_\lambda$ symmetric in $(\mu,\lambda)$ and $\chi_\nu\chi_\sigma$ in $(\nu,\sigma)$, and $\Gamma$ is Hermitian,
$D^{\rm raw}_{\lambda\sigma\mu\nu}=\overline{D^{\rm raw}_{\mu\nu\lambda\sigma}}$, so the average over the four permutations
$D^S=\tfrac14\sum_{\rm perm}D^{\rm raw}=\tfrac12\operatorname{Re}\big(D^{\rm raw}_{\mu\nu\lambda\sigma}+D^{\rm raw}_{\lambda\nu\mu\sigma}\big)$
is exactly real: the imaginary parts cancel by complex conjugation, and only $D^S$ is written (the largest imaginary part of
the symmetrized element is printed as a check). The 1-RDM is $\gamma^S_{\mu\lambda}=\tfrac12\operatorname{Re}\big(\gamma_{\mu\lambda}+\gamma_{\lambda\mu}\big)$,
$\gamma_{\mu\lambda}=\sum_c\sum_{pr}\overline{C^c_{\mu p}}\bar\gamma_{pr}C^c_{\lambda r}$.

\paragraph{Files.}
\texttt{.basis}: a header line, then one line \texttt{Z x y z exponent nx ny nz} per primitive, Large primitives first.
\texttt{.dm2}: Fortran-unformatted-style records (\texttt{int32} record length, four \texttt{int32} indices $p,q,r,s$,
\texttt{float64} value, \texttt{int32} record length), 1-based primitive indices, terminated by a record of four zero indices.
Both $D_{pqrs}$ and its electron-exchange partner $D_{qpsr}$ are present, with $r\le p$ and $s\le q$ (the
\texttt{RHO2\_OPS} convention without \texttt{\$SYMRED}); the stored value is the sum of the symmetry-equivalent elements,
$m_{pr}m_{qs}D^S_{pqrs}$ with $m=2$ when the two indices differ and $1$ otherwise.
\texttt{.dm1}: records (\texttt{int32} length, \texttt{int32} $\mu$, \texttt{int32} $\lambda$, \texttt{float64}
$\gamma^S_{\mu\lambda}$, \texttt{int32} length) for $\mu\ge\lambda$, terminated by $\mu=\lambda=0$.

\paragraph{Algorithm.} $\Gamma$ is never held as an $n^4$ array: it is generated one slab $\Gamma_{p\,\cdot\cdot\cdot}$ at a time
from the $d$-blocks. For a block of rows $\mu$ of one sector, a sweep over $p$ accumulates
$R^c_{\mu qrs}=\sum_p\overline{C^c_{\mu p}}\Gamma_{pqrs}$ and $W^c_{\mu pqs}=\sum_rC^c_{\mu r}\Gamma_{pqrs}$ (\texttt{zgemm}); then, for each $\mu$,
$Y_{\mu\lambda;qs}=\tfrac12\sum_c\big(\sum_rC^c_{\lambda r}R^c_{\mu qrs}+\sum_p\overline{C^c_{\lambda p}}W^c_{\mu pqs}\big)$ for $\lambda\le\mu$ (the
two-term average already is the permutation average over $\mu\leftrightarrow\lambda$), and finally
$E_{\mu\lambda;\nu\sigma}=\sum_{c_2}\sum_{qs}\overline{C^{c_2}_{\nu q}}Y_{\mu\lambda;qs}C^{c_2}_{\sigma s}$, whose real part is $D^S$.
Only $\nu\le\mu$ is computed and the electron-exchange partner is written alongside; each record is written to disk at once, so the
memory is $2N_c B n^3$ complex numbers for a block of $B$ rows (\texttt{CC\_RDM\_MAXMEM} sets $B$; more rows per block means
fewer sweeps over $\Gamma$). The cost is about $4P^4n$ flops for $P$ primitives and $n$ spinors, plus the sweeps over $\Gamma$.

\paragraph{Checks and validation.} Printed after each run: $\sum\gamma^S S=N$ and $\sum D^SSS=N(N-1)$ with the primitive
overlaps. With \texttt{RERDMFT\_CC\_RDM\_GAU\_CHECK=1}, $V_{ee}=\tfrac12\sum D\,(\mu\lambda|\nu\sigma)$ is also recomputed with the
integrals of the primitives themselves (small systems): LiH/6-31G (NON\_REL, X2C and 4-component, with the Small primitives) and
water/cc-pVDZ with a frozen core agree with the molecular-basis two-body energy to $10^{-11}$--$10^{-13}$~Eh with dense integrals
(with RI the difference is the RI error). \texttt{RHO2\_OPS} reads the files directly and reports $\mathrm{Tr}[{}^2D]=N(N-1)/2$ (for LiH,
$V_{ee}$ from the intracule converges to the molecular value as its radial/angular quadratures are refined).

\subsection{Numerical check}
\label{sec:lambda-check}

PySCF~2.14 (\texttt{GCCSD}) is the reference. All checks use complex integrals: a Hartree--Fock reference of LiH is rotated
inside the occupied and virtual spaces by random unitaries $e^{A}$, $A=-A^\dagger$ complex (and, in a second set, with a
small complex occupied--virtual mixing so that $f_{ov}\neq0$ and $E_c$ acquires a small imaginary part).
\begin{itemize}
\item The equations above, evaluated from the same data, reproduce PySCF's intermediates and one full $\Lambda$ update
 to $10^{-15}$, and the $d$-blocks and $\mathcal R$ array to $10^{-16}$.
\item $\partial\mathcal L/\partial t_\mu$ at PySCF's converged $(t,\lambda)$ vanishes to $3\times10^{-12}$ for all singles and
 all antisymmetry-unique doubles.
\item Central finite differences of $\operatorname{Re}E_{\rm CC}$ (amplitudes re-solved) with respect to a Hermitian
 perturbation of $h$ and of the antisymmetrized integrals equal $\sum h'_{pq}\bar\gamma_{pq}$ and
 $\tfrac14\sum\anti{p}{q}{r}{s}'\Gamma_{pqrs}$ to $10^{-8}$.
\item $\sum h\bar\gamma+\tfrac14\sum\anti{p}{q}{r}{s}\Gamma_{pqrs}=\operatorname{Re}(E_0+E_c)$ to $10^{-14}$.
\end{itemize}

\section{Implementation (\texttt{src/CC/Uccsd.\{h,cpp\}})}
\label{sec:impl}

The code is a transcription of Section~4 with two reorganizations that do not change the equations. Every
contraction is a two-operand einsum evaluated as one \texttt{zgemm}.

\begin{itemize}
\item \textbf{No stored $\langle ma\Vert ef\rangle$ or $\langle ab\Vert ef\rangle$ blocks.} The integrals are
 supplied as a Coulomb factorization $(pq|rs)=\sum_P \mathrm{bra}_P(p,q)\,\mathrm{ket}_P(r,s)$ (RI: bra = ket = $B_P$;
 Cholesky: $\mathrm{bra}_L(a,c)=V_L(c,a)$, $\mathrm{ket}_L(b,d)=V_L^*(b,d)$), and one slab with a fixed first index,
 $\langle p\,a|e f\rangle$, is rebuilt by a single \texttt{zgemm} whenever it is needed. A loop over occupied $m$
 accumulates every term containing $\langle ma\Vert ef\rangle$ ($F_{ae}$, the $T_1$ equation, $W_{mbej}$, the
 $P(ij)\,t_i^e\langle ab\Vert ej\rangle$ term and the intermediate $Z$ below); a loop over virtual $a$ applies
 the particle--particle ladder. Memory is $O(o^2v^2+N_{\rm aux}n^2)$, plus $O(N_{\rm aux}v^2)$ for the cached
 second factor. For dense integrals the same slabs are read element by element.
\item \textbf{The ladder is applied as $\sum_{ef}\tau_{ij}^{ef}\langle ab|ef\rangle$.} Because $\tau$ is antisymmetric
 in $(ef)$, $\tfrac12\sum_{ef}\tau_{ij}^{ef}\anti{a}{b}{e}{f}=\sum_{ef}\tau_{ij}^{ef}\dir{a}{b}{e}{f}$, so no
 antisymmetrized $vvvv$ slab is formed. The two other pieces of $\tfrac12\sum\tau W_{abef}$ are
 $-P(ab)\sum_m t_m^b Z_{ijm}^a$ with $Z_{ijm}^a=\tfrac12\sum_{ef}\tau_{ij}^{ef}\anti{a}{m}{e}{f}$, and the
 $\tau\tau$ term, which is moved into $W_{mnij}$ by doubling its coefficient there
 ($\tfrac14\to\tfrac12$): $\tfrac12\cdot\tfrac14+\tfrac12\cdot\tfrac14=\tfrac12\cdot\tfrac12$.
\end{itemize}

Amplitudes are updated by Jacobi sweeps, $t\leftarrow(\mathrm{RHS}-\sigma t)/(D-\sigma)$ with level shift $\sigma$,
accelerated by DIIS on the antisymmetry-unique amplitudes (real weights, \texttt{Utils/DIIS.h}). The Fock matrix
is built from the integrals (frozen occupied spinors included), so $f_{ov}\neq0$ and non-canonical orbitals are
handled. Checked against PySCF 2.14: LiH/6-31G, water/cc-pVDZ with frozen core and CO/cc-pVDZ with frozen core
(dense and Cholesky agree to $10^{-10}$--$10^{-8}$~Eh with PySCF's RCCSD), and the $c\to\infty$ limit of the
X2C and 4-component references reproduces the non-relativistic result.

\end{document}
